\begin{frame}{Amplitudes medidas em bancada}
  \small
  \renewcommand{\arraystretch}{0.93}
  \begin{columns}[T,onlytextwidth]
    \begin{column}{0.48\textwidth}
      \centering
      \textbf{Circuito 1: $V_i=78$--$155\,\mathrm{mV}$}

      \vspace{0.10cm}
      \begin{tabular}{rr}
        \toprule
        $f$ (Hz) & $V_o$ (V) \\
        \midrule
        0,6 & 0,750 \\
        1,7 & 1,050 \\
        2,5 & 0,960 \\
        4,3 & 0,730 \\
        5,2 & 0,640 \\
        6,0 & 0,570 \\
        7,7 & 0,460 \\
        10,8 & 0,340 \\
        17,2 & 0,210 \\
        25,8 & 0,140 \\
        43,0 & 0,060 \\
        \bottomrule
      \end{tabular}
    \end{column}
    \begin{column}{0.48\textwidth}
      \centering
      \textbf{Circuito 2: $V_i=4{,}7$--$5{,}1\,\mathrm{V}$}

      \vspace{0.10cm}
      \begin{tabular}{rr}
        \toprule
        $f$ (Hz) & $V_o$ (V) \\
        \midrule
        2,6 & 4,700 \\
        7,0 & 5,000 \\
        10,5 & 5,100 \\
        17,5 & 5,000 \\
        21,0 & 4,700 \\
        24,5 & 4,600 \\
        31,0 & 3,900 \\
        43,7 & 2,700 \\
        70,0 & 1,300 \\
        105,0 & 0,600 \\
        175,0 & 0,220 \\
        \bottomrule
      \end{tabular}
    \end{column}
  \end{columns}

  \vspace{0.14cm}
  \centering\scriptsize
  Amplitudes de saída convertidas para volts.
  A tensão de entrada foi medida em cada frequência da varredura.
\end{frame}
